\chapter{Introduction}
\label{ch:introduction}

This chapter states the research context and problem of SENTINEL---a cognitive two-tier architecture automating threat detection and incident response in Security Operations Centres (SOCs)---with the objectives, research questions, scope, and contributions of the thesis.

\section{Background and Motivation}

A mid-sized Security Operations Centre now ingests millions of log events every day~\cite{siem2021}, almost all recording ordinary activity. Because missing an intrusion costs far more than examining a harmless event, monitoring is deliberately tuned toward sensitivity, and practitioners consequently report that only a small fraction of what reaches them merits any action~\cite{alertfatigue2022, alahmadi2022}. Every additional sensor raises the alert count while the number of specialists able to read those alerts stays fixed, so the binding constraint of a modern SOC is not detection but the allocation of scarce human attention, which cannot be multiplied by buying more servers.

Three layers of countermeasure stand between the stream and the analyst: firewalls and intrusion detection systems matching static signatures, SIEM platforms correlating centralised telemetry with hand-written rules, and SOAR platforms executing pre-programmed playbooks~\cite{soar2025}. Whatever they cannot settle arrives on an analyst's queue, where the work becomes manual---retrieving the raw log, reconstructing what the source address did around the flagged moment, checking the destination against the asset inventory, and identifying the technique the behaviour resembles before deciding to block, dismiss, or escalate. Every step is reading and cross-referencing, and reading does not parallelise.

Two weaknesses follow, and they are the reasons this thesis exists. The first is time: an analyst needs minutes per alert while the stream arrives continuously, so a queue whose arrival rate exceeds its service rate is cleared by abandonment rather than completion. The second is context: the evidence deciding a case is scattered across records and buried in unstructured fields, so what is needed is not a pattern to match but a judgement to form. Static rules cannot express that, and signatures and playbooks recognise only what has already been described. Large Language Models supply the missing capability at three new costs: each act of reasoning takes seconds while data keeps arriving~\cite{vaswani2017}; the reasoning layer becomes a target, since a log field can carry wording an attacker composed to steer the model~\cite{promptinjection2023} and any conclusion must be provably unaltered to serve as evidence; and internal logs may not leave the organisation, so the model must run on site on finite hardware~\cite{nist80207}.

The founding observation is the cost gap between the two kinds of decision: a statistical rule decides in under a millisecond, a language model in seconds. Four to five orders of magnitude make routing every event through the most expensive tier a systematic waste. The objective is therefore to design, build, and empirically evaluate SENTINEL, a two-tier architecture built on \textit{placing cheap decisions before expensive ones}: a fast filtering tier settling most traffic outright, and a reasoning tier called upon only for what remains, behind protective guardrails and a cryptographically sealed audit trail, entirely on local infrastructure. What makes this a research problem rather than an engineering one is that each of those three costs must be \textit{measured}, on real data, with metrics that cannot flatter themselves.

\section{Objectives, Research Questions, Scope and Contributions}
\label{sec:rq}

The general objective is operationalised as three research tasks, each paired one-to-one with a research question and the contribution it is intended to produce. Chapter~\ref{ch:experiments_evaluation} reports results under the same three headings.

The first task is filtration economics: a stateless dual-stage filtration pipeline operating at wire-speed, resolving the majority of telemetry at the cheap tier without consuming GPU resources. RQ1 asks: \textit{How can a multi-tiered architecture be designed to autonomously offload security telemetry at wire-speed, thereby addressing the latency bottlenecks and resource consumption inherent in Large Language Models?} The intended contribution is to quantify the principle of placing cheap decisions before expensive ones, and to demonstrate by measurement that the offload rate is a function of the attack base rate of the stream rather than a constant of the system---a property not stated explicitly in the related literature.

The second task is the security of the reasoning layer: to protect the language model structurally rather than probabilistically, and to seal the forensic record so that any alteration is detectable. RQ2 asks: \textit{What methodologies and security mechanisms can effectively withstand adversarial risks (such as log-substrate prompt injection) to protect the AI's reasoning flow, while simultaneously ensuring the integrity and tamper-evidence of forensic audit trails?} The intended contribution is a defence based on Delimited Data Encapsulation with random cryptographic nonces, whose essential property is independence from the block rate of any classifier: encapsulated content is processed as passive text regardless of what it says, unlike probabilistic semantic filters whose effectiveness degrades against novel variants~\cite{owasp_llm2025}.

The third task covers reasoning, attribution and human workload: a stateful reasoning engine that retrieves domain knowledge, cites the evidence it used, and routes to a human the cases that warrant one. RQ3 asks: \textit{How can stateful reasoning capabilities and domain-specific knowledge retrieval be integrated into an AI Agent to automate intrusion analysis, attribute attack techniques, and generate transparent forensic reports, thereby reducing and accurately triaging the workload requiring human analysts?} The intended contribution is to replace rigid SOAR playbooks with a stateful agent over a STIX-aligned MITRE ATT\&CK knowledge base, coupled with a deterministic RAG-anchoring guardrail requiring every emitted technique identifier to be present in that batch's retrieved documents---turning hallucination control from a hope placed in the model into a verifiable architectural constraint.

The scope is bounded as follows. The object of study is the SENTINEL architecture: a Tier 1 stateless filter, a Tier 1.75 semantic cache, a Tier 2 stateful agent running a locally hosted quantised model, hybrid retrieval over MITRE ATT\&CK knowledge, and the cryptographic controls protecting both. Evaluation rests on two primary datasets, CSE-CIC-IDS2018~\cite{ids2018} for wide-area network telemetry and CSIC 2010~\cite{csic2010} for HTTP Layer 7 attacks; DAPT2020~\cite{dapt2020} for multi-stage chains and a synthetic zero-day set serve as proof-of-concept evidence only and support no contribution claim. Adversarial evaluation uses 823 samples partitioned by provenance and never merged: 703 from published corpora (AdvBench~\cite{advbench2023}, deepset~\cite{deepset2023}, jackhhao~\cite{jackhhao2023}) as the headline denominator, 80 author-composed, and 40 cross-check. The system runs on a single air-gapped workstation; automated responses are restricted to blocking IP addresses, logging, or escalating to an analyst.

\section{Structure of the Thesis}

Chapter~\ref{ch:theoretical_background} systematises the theoretical foundations and surveys related work to establish the research gap. Chapter~\ref{ch:system_design_implementation} states the research methodology and the technical design of SENTINEL. Chapter~\ref{ch:experiments_evaluation} presents the quantitative results, organised to follow the three research questions above, and Chapter~\ref{ch:conclusion} concludes.
